% ======================================================================
%  sec07.tex —— 第 7 节 映入仿射概形的态射
%  底本：原书 pp.42–43
%  原书标题：7. Morphisms to affine schemes
%            7.1 A criterion for being affine / 7.2 Applications
% ======================================================================

\section{映入仿射概形的态射}\label{sec:7}

\subsection{仿射性的判别准则}\label{sec:7.1}

\begin{construction}\label{con:7.1.1}
设 $X$ 与 $Y$ 为概形。对任意态射 $f:X\to Y$，按定义有态射 $\alpha:\cO_Y\to f_*\cO_X$。取整体截面，得到同态
\[
\alpha_Y:\cO_Y(Y)\to(f_*\cO_X)(Y)\simeq\cO_X(X).
\]
可以验证，这给出函子
\[
\Sch\to\CRing^{\mathrm{op}},\qquad X\mapsto\cO_X(X),
\]
它把态射 $f:X\to Y$ 送到上述同态 $\alpha_Y$。
\end{construction}

\begin{yizhu}
构造 \ref{con:7.1.1} 的函子是反变的：态射 $f:X\to Y$ 被送到 $f$ 在整体截面上的拉回 $\alpha_Y$，方向由 $\cO_Y(Y)$ 指向 $\cO_X(X)$。这一步之所以合法，是因为整体截面是层的函子性中最容易写出的部分；真正有内容的是下一个定理，它说对仿射的 $Y$，这个显然的映射竟是双射。
[依据：原书 p.42 第 1 段与定理 7.1.2]
\end{yizhu}

\begin{theorem}\label{thm:7.1.2}
概形 $Y$ 仿射，当且仅当对任意概形 $X$，自然映射
\begin{equation}\tag{7.1}\label{eq:7.1}
\Hom_{\Sch}(X,Y)\to\Hom_{\CRing}\bigl(\cO_Y(Y),\cO_X(X)\bigr)
\end{equation}
是双射。
\end{theorem}

\begin{proof}
先证明「仅当」方向。设 $Y\simeq\Spec(R)$ 为仿射概形，要证自然映射
\begin{equation}\tag{7.2}\label{eq:7.2}
\Hom_{\Sch}(X,\Spec(R))\to\Hom_{\CRing}\bigl(R,\cO_X(X)\bigr)
\end{equation}
是双射。

按如下方式构造反方向的映射。设 $h:R\to\cO_X(X)$ 为同态。对任意仿射开子概形 $U\subseteq X$，由命题--定义 \ref{prop:5.1.9}，映射
\[
\Hom_{\Sch}(U,\Spec(R))\to\Hom_{\CRing}\bigl(R,\cO_X(U)\bigr)
\]
是双射。令 $f_U:U\to\Spec(R)$ 为对应于合成
\[
R\xrightarrow{\ h\ }\cO_X(X)\xrightarrow{\ (-)|_U\ }\cO_X(U)
\]
的态射。展开构造可以验证：对仿射开子概形 $U\subseteq V\subseteq X$，有 $f_U=f_V|_U$。仿射开子集构成 $X$ 的拓扑基。利用习题 \ref{exc:1.2.7} 可知，存在唯一态射 $p:X\to\Spec(R)$，使对任意仿射开子概形 $U\subseteq X$ 都有 $f_U=p|_U$。上述构造 $h\mapsto p$ 给出映射
\begin{equation}\tag{7.3}\label{eq:7.3}
\Hom_{\CRing}\bigl(R,\cO_X(X)\bigr)\to\Hom_{\Sch}(X,\Spec(R)).
\end{equation}

可以验证 (7.2) 与 (7.3) 互为逆映射。具体地说，$\eqref{eq:7.3}\circ\eqref{eq:7.2}=\id$ 由态射 $p$ 的唯一性得出；而 $\eqref{eq:7.2}\circ\eqref{eq:7.3}=\id$ 由下述事实得出：$\cO_X(X)$ 中的元素被它在各个 $\cO_X(U)$ 中的限制决定。

现在由「仅当」方向推出「若」方向。设概形 $Y$ 使 (7.1) 对任意概形 $X$ 都是双射。记 $R:=\cO_Y(Y)$。由已证的「仅当」方向，有
\[
\Hom_{\Sch}(Y,\Spec(R))\to\Hom_{\CRing}(R,R)
\]
是双射。特别地，$\id_R$ 对应典范态射 $q:Y\to\Spec(R)$。

此外，按构造，对任意概形 $X$，下图交换：
\[
\begin{tikzcd}[column sep=large,row sep=large]
\Hom_{\Sch}(X,Y)\arrow[r,"\simeq"]\arrow[d,"q\circ-"'] & \Hom_{\CRing}\bigl(R,\cO_X(X)\bigr)\arrow[d,equal]\\
\Hom_{\Sch}(X,\Spec(R))\arrow[r,"\simeq"'] & \Hom_{\CRing}\bigl(R,\cO_X(X)\bigr)
\end{tikzcd}
\]
其中水平映射分别是把 (7.1) 用于 $Y$ 与 $\Spec(R)$，而它们或是按假设、或是按已证的「仅当」方向而为双射。于是与 $q$ 的合成诱导双射
\[
\Hom_{\Sch}(X,Y)\xrightarrow{\ \simeq\ }\Hom_{\Sch}(X,\Spec(R))
\]
对任意 $X\in\Sch$ 成立。由 Yoneda 引理，$q$ 是同构，因此 $Y$ 仿射。
\end{proof}

\begin{yizhu}
定理 \ref{thm:7.1.2} 在 Stacks 中有逐字对应的引理：对仿射的 $Y$，自然映射
\[
\Mor(X,Y)\to\Hom_{\CRing}\bigl(\Gamma(Y,\cO_Y),\Gamma(X,\cO_X)\bigr)
\]
是双射，并由此得到 $\Gamma$ 与 $\Spec$ 的伴随关系。区别只在于 Stacks 先建立仿射情形的等价，再把它作为推论使用。
[来源：Stacks Project, Tag 01I1（Lemma 26.6.4, The category of affine schemes），\url{https://stacks.math.columbia.edu/tag/01I1}（2026-11-17 访问）。]
\end{yizhu}

\subsection{应用}\label{sec:7.2}

\begin{corollary}\label{cor:7.2.1}
对任意概形 $X$，存在唯一态射 $X\to\Spec(\Z)$。换言之，$\Spec(\Z)\in\Sch$ 是终对象。
\end{corollary}

\begin{yizhu}
推论 \ref{cor:7.2.1} 的对偶陈述是：空概形是 $\Sch$ 的初始对象。两者合起来说明 $\Sch$ 没有零对象——它既有初始对象又有终对象，而二者并不同构。这一点与 $\CRing$ 的情形形成对照：在 $\CRing$ 中零环既是初始对象又是终对象。
[译者按（上下文推断，无外部依据）：由终对象与初始对象的定义直接比较。]
\end{yizhu}

\begin{corollary}\label{cor:7.2.2}
对任意概形 $X$，存在唯一态射
\[
q:X\to\Spec\bigl(\cO_X(X)\bigr),
\]
它在整体截面之间诱导恒等同态。
\end{corollary}

\begin{corollary}\label{cor:7.2.3}
嵌入函子 $\mathsf{Aff}\to\Sch$ 有典范左伴随
\[
\Sch\to\mathsf{Aff},\qquad X\mapsto\Spec\bigl(\cO_X(X)\bigr).
\]
\end{corollary}

\begin{yizhu}
推论 \ref{cor:7.2.3} 中的左伴随通常称为\emph{仿射化}：它把任意概形 $X$ 送到仿射概形 $\Spec(\cO_X(X))$，而推论 \ref{cor:7.2.2} 的 $q$ 正是这一伴随的单位。对仿射 $X=\Spec(R)$，单位是同构 $R\simeq\cO_X(X)$，故仿射概形是此伴随的不动点。
[依据：原书 p.43 推论 7.2.2 与 7.2.3]
\end{yizhu}

\begin{exercise}\label{exc:7.2.4}
证明嵌入函子 $\mathsf{Aff}\to\Sch$ 没有右伴随。
\end{exercise}

